% Transcripción LaTeX reconstruida a partir de certamen-1.pdf.
% La fuente original del PDF no se encontró entre los archivos disponibles.
\documentclass[11pt,a4paper]{article}
\usepackage[utf8]{inputenc}
\usepackage[T1]{fontenc}
\usepackage[spanish,es-nodecimaldot]{babel}
\usepackage{amsmath,amssymb,amsthm}
\usepackage{geometry}
\usepackage{lmodern}
\geometry{margin=2.5cm}

\newcommand{\R}{\mathbb{R}}
\newcommand{\Z}{\mathbb{Z}}
\newcommand{\N}{\mathbb{N}}
\newcommand{\set}[1]{\left\{#1\right\}}
\newcommand{\opm}{\mathbin{\oplus_m}}

\title{Certamen 1 --- Teoría de Grupos}
\author{José Ignacio Rosas Sepúlveda}
\date{28 de julio de 2026}

\begin{document}
\maketitle

\section*{Problema 1}

Sea $m$ cualquier entero positivo fijo y sea
\[
S=\set{0,1,2,\ldots,m-1}.
\]
Definir en $S$ la suma reducida al conjunto $S$ y probar que el sistema obtenido es un grupo de orden $m$.

\subsection*{Solución}

Sean $a,b\in S$. Como $0\leq a+b\leq 2m-2$, para reducir la suma $a+b$ al conjunto $S$ basta restar $m$ cuando $a+b\geq m$. Definimos entonces
\[
a\opm b:=
\begin{cases}
a+b, & 0\leq a+b<m,\\
a+b-m, & m\leq a+b\leq 2m-2.
\end{cases}
\]

\subsubsection*{A1) Clausura}

Sean $a,b\in S$. Si $a+b<m$, entonces
\[
0\leq a\opm b=a+b\leq m-1,
\]
y, por tanto, $a\opm b\in S$. Si $a+b\geq m$, entonces
\[
0\leq a\opm b=a+b-m\leq m-2,
\]
y nuevamente $a\opm b\in S$. Por lo tanto,
\[
a,b\in S\quad\Longrightarrow\quad a\opm b\in S.
\]

\subsubsection*{A2) Asociatividad}

Sean $a,b,c\in S$ y denotemos $T=a+b+c$. Como $0\leq a,b,c\leq m-1$, se tiene $0\leq T\leq 3m-3$. Compararemos las dos formas de asociar la operación.

\paragraph{Caso 1: $0\leq T<m$.}
Como $a+b\leq T<m$ y $b+c\leq T<m$, ninguna de las dos sumas parciales requiere reducción. Luego
\[
(a\opm b)\opm c=(a+b)+c=T=a+(b+c)=a\opm(b\opm c).
\]

\paragraph{Caso 2: $m\leq T<2m$.}
Para el miembro izquierdo hay dos posibilidades. Si $a+b<m$, entonces
\[
(a\opm b)\opm c=(a+b)\opm c=T-m,
\]
pues $m\leq T<2m$. Si $a+b\geq m$, entonces $a\opm b=a+b-m$ y, como $0\leq T-m<m$,
\[
(a\opm b)\opm c=(a+b-m)+c=T-m.
\]
Por consiguiente, $(a\opm b)\opm c=T-m$. Aplicando el mismo razonamiento a la suma parcial $b+c$, se obtiene $a\opm(b\opm c)=T-m$. Por tanto,
\[
(a\opm b)\opm c=a\opm(b\opm c).
\]

\paragraph{Caso 3: $2m\leq T\leq 3m-3$.}
Como $c\leq m-1$,
\[
a+b=T-c\geq 2m-(m-1)=m+1,
\]
de modo que $a\opm b=a+b-m$. Además, $m\leq T-m\leq 2m-3$, por lo que la segunda suma también debe reducirse:
\[
(a\opm b)\opm c=(a+b-m)\opm c=T-2m.
\]
Análogamente, como $a\leq m-1$, se tiene $b+c\geq m+1$ y, por tanto,
\[
a\opm(b\opm c)=T-2m.
\]
Luego $(a\opm b)\opm c=a\opm(b\opm c)$. Los tres casos posibles para $T$ muestran que $\oplus_m$ es asociativa.

\subsubsection*{A3) Existencia de identidad}

El candidato a elemento identidad es $0\in S$. Sea $a\in S$. Como $a<m$,
\[
a\opm 0=a+0=a,\qquad 0\opm a=0+a=a.
\]
Por lo tanto, $a\opm 0=0\opm a=a$ para todo $a\in S$. Luego el elemento identidad es $0$.

\subsubsection*{A4) Existencia de inverso}

Sea $a\in S$. Si $a=0$, su inverso es $0$, pues $0\opm 0=0$. Si $a\neq 0$, entonces $1\leq a\leq m-1$. Definimos
\[
a^{-1}:=m-a.
\]
Como $1\leq m-a\leq m-1$, se tiene $a^{-1}\in S$. Además,
\[
a+a^{-1}=a+(m-a)=m,
\]
por lo que, de acuerdo con la segunda parte de la definición, $a\opm a^{-1}=m-m=0$. Como la suma ordinaria de enteros es conmutativa, también $a^{-1}\opm a=0$. Luego todo elemento de $S$ posee inverso en $S$.

\subsection*{Orden del grupo}

El conjunto $S=\set{0,1,2,\ldots,m-1}$ contiene exactamente $m$ elementos. Por tanto, $|S|=m$.

\subsection*{Resultado}

La operación $\oplus_m$ satisface clausura, asociatividad, existencia de identidad y existencia de inverso. En consecuencia,
\[
(S,\oplus_m)
\]
es un grupo de orden $m$.

\newpage
\section*{Problema 2}

Sea el grupo $(\R,+)$. Muestre que el grupo $(\Z,+)$ es un subgrupo normal de $(\R,+)$, encuentre el grupo cociente y determine a qué grupo es isomorfo.

\subsection*{Solución}

\subsubsection*{$(\Z,+)$ es subgrupo de $(\R,+)$}

Se tiene $\Z\subset\R$ y $0\in\Z$. Sean $m,n\in\Z$. Entonces $m-n\in\Z$. Por el criterio de subgrupo,
\[
(\Z,+)\leq(\R,+).
\]

\subsubsection*{Normalidad}

El grupo $(\R,+)$ es abeliano. Por lo tanto, todo subgrupo de $(\R,+)$ es normal. En particular,
\[
(\Z,+)\trianglelefteq(\R,+).
\]
De forma explícita, para todo $x\in\R$,
\[
x+\Z-x=\set{x+n-x:n\in\Z}=\Z.
\]

\subsubsection*{Grupo cociente}

Los elementos de $\R/\Z$ son los cosets
\[
x+\Z=\set{x+n:n\in\Z},\qquad x\in\R.
\]
La operación del grupo cociente es la suma de cosets,
\[
(x+\Z)+(y+\Z):=(x+y)+\Z.
\]
Como $\Z$ es normal en $\R$, esta operación está bien definida.

Todo coset posee un representante en el intervalo $[0,1)$. En efecto, para $x\in\R$ tomamos $n=\lfloor x\rfloor\in\Z$ y $t=x-n\in[0,1)$. Entonces $x+\Z=t+\Z$. Además, este representante es único: si $s,t\in[0,1)$ y $s+\Z=t+\Z$, entonces $s-t\in\Z$; pero $-1<s-t<1$, de modo que $s-t=0$ y $s=t$.

Por lo tanto, el grupo cociente puede escribirse como
\[
\R/\Z=\set{t+\Z:0\leq t<1},
\]
con la operación
\[
(s+\Z)+(t+\Z)=(s+t)+\Z.
\]

\subsubsection*{Isomorfismo con el grupo circular}

Consideremos el grupo circular
\[
(S^1,\cdot),\qquad S^1=\set{z\in\mathbb{C}:|z|=1}.
\]
Definimos la aplicación
\[
\varphi:(\R,+)\longrightarrow(S^1,\cdot),\qquad \varphi(x)=e^{2\pi i x}.
\]

\paragraph{Homomorfismo.} Sean $x,y\in\R$. Entonces
\[
\varphi(x+y)=e^{2\pi i(x+y)}=e^{2\pi i x}e^{2\pi i y}=\varphi(x)\varphi(y).
\]
Por lo tanto, $\varphi$ es un homomorfismo.

\paragraph{Núcleo.} Se tiene
\[
\ker\varphi=\set{x\in\R:\varphi(x)=1}
=\set{x\in\R:e^{2\pi i x}=1}=\Z.
\]

\paragraph{Sobreyectividad.} Sea $z\in S^1$. Existe $\theta\in[0,2\pi)$ tal que $z=e^{i\theta}$. Tomando $x=\theta/(2\pi)\in\R$, se obtiene
\[
\varphi(x)=e^{2\pi i(\theta/(2\pi))}=e^{i\theta}=z.
\]
Por lo tanto, $\operatorname{Im}\varphi=S^1$.

Por el teorema fundamental de los homomorfismos,
\[
\R/\ker\varphi\cong\operatorname{Im}\varphi.
\]
Como $\ker\varphi=\Z$ y $\operatorname{Im}\varphi=S^1$, concluimos que
\[
\R/\Z\cong(S^1,\cdot).
\]

\subsection*{Resultado}

Se ha demostrado que
\[
(\Z,+)\trianglelefteq(\R,+).
\]
El grupo cociente está formado por los cosets $x+\Z$ y su operación es
\[
(x+\Z)+(y+\Z)=(x+y)+\Z.
\]
Cada coset tiene un único representante en $[0,1)$ y, finalmente,
\[
\R/\Z\cong(S^1,\cdot).
\]

\end{document}
